\section{2020 Algebra \& Number Theory}

\subsection{Individual \& Team}
\begin{exercise}
\label{algebra:2020:1}
Let $F$ be a field of characteristic zero. Consider the polynomial ring $F[x_1, \ldots, x_n]$.

\begin{enumerate}
\item Prove Newton's identity over the field $F$:
\[p_k - p_{k-1}e_1 + \cdots + (-1)^{k-1}p_1e_{k-1} + (-1)^k k e_k = 0,\]
where
\[e_k(x_1, \ldots, x_n) = \sum_{1 \leq i_1 < \cdots < i_k \leq n} x_{i_1} \cdots x_{i_k}\]
for $1 \leq k \leq n$, $e_0 = 1$, $e_k = 0$ when $k > n$, and
\[p_k(x_1, \ldots, x_n) = x_1^k + \cdots + x_n^k.\]

\item Prove that over the field $F$ of characteristic zero, an $n \times n$ matrix $A$ is nilpotent if and only if the trace of $A^k$ is equal to zero for all $k = 1, 2, \ldots$.

Hint: use Part (a).

\item Prove that over the field $F$ of characteristic zero, two $n \times n$ matrices $A$ and $B$ have the same characteristic polynomial if and only if the trace of $A^k$ and trace of $B^k$ are equal for all $k = 1, 2, \ldots$.

Hint: use Part (a).
\end{enumerate}
\end{exercise}

\begin{solution}
\subsection*{Part (1): Newton's Identity}

\textbf{Prerequisites:}
To follow this proof, you should be familiar with \textit{elementary symmetric polynomials} ($e_k$), \textit{power sum symmetric polynomials} ($p_k$), and the use of \textit{generating functions} (a way to encode a sequence of numbers as coefficients of a polynomial or power series).

\textbf{Steps:}
\begin{enumerate}
    \item Define the generating function for the elementary symmetric polynomials:
    \[ E(t) = \sum_{k=0}^n e_k t^k = \prod_{i=1}^n (1 + x_i t). \]
    \item Take the logarithmic derivative of $E(t)$. By the chain rule, $(\ln E(t))' = E'(t)/E(t)$. Since $E(t)$ is a product, its logarithm is a sum:
    \[ \frac{E'(t)}{E(t)} = \frac{d}{dt} \sum_{i=1}^n \ln(1 + x_i t) = \sum_{i=1}^n \frac{x_i}{1 + x_i t}. \]
    \item Expand each term in the sum as a geometric series in $t$:
    \[ \sum_{i=1}^n x_i \sum_{j=0}^\infty (-x_i t)^j = \sum_{i=1}^n \sum_{k=1}^\infty (-1)^{k-1} x_i^k t^{k-1} = \sum_{k=1}^\infty (-1)^{k-1} p_k t^{k-1}. \]
    \item Multiply both sides by $E(t)$ to clear the denominator:
    \[ E'(t) = E(t) \cdot \left( \sum_{m=1}^\infty (-1)^{m-1} p_m t^{m-1} \right). \]
    \item Match the coefficients of $t^{k-1}$ on both sides. On the left, the coefficient of $t^{k-1}$ in $E'(t) = \sum j e_j t^{j-1}$ is $k e_k$. On the right, use the Cauchy product of the two series.
\end{enumerate}

\textbf{Intuition:}
Generating functions act like "suitcases" that hold all the symmetric polynomials at once. By applying calculus (differentiation and logarithms) to these suitcases, we can uncover hidden relationships between the different types of symmetric polynomials. Newton's identity is essentially a recursive way to build complex "power sums" from simpler "building blocks" (the elementary symmetric polynomials).

\par\vspace{1.5em}

\subsection*{Part (2): Nilpotent Matrices}

\textbf{Prerequisites:} 
Recall that a matrix $A$ is \textit{nilpotent} if $A^m = 0$ for some $m$. This is equivalent to all its eigenvalues being zero. The \textit{trace} of $A^k$ is the sum of the $k$-th powers of its eigenvalues.

\textbf{Steps:}
\begin{enumerate}
    \item Let $\lambda_1, \dots, \lambda_n$ be the eigenvalues of $A$. Then $\mathrm{Tr}(A^k) = p_k(\lambda_1, \dots, \lambda_n)$.
    \item If $\mathrm{Tr}(A^k) = 0$ for all $k$, then $p_k = 0$ for all $k$.
    \item Using Newton's identity from Part (1), if $p_1 = p_2 = \dots = p_n = 0$, then $e_1 = e_2 = \dots = e_n = 0$.
    \item The characteristic polynomial of $A$ is $\chi_A(t) = t^n - e_1 t^{n-1} + \dots + (-1)^n e_n$. With all $e_k = 0$, $\chi_A(t) = t^n$.
    \item By the Cayley-Hamilton theorem, $A^n = 0$, meaning $A$ is nilpotent.
\end{enumerate}

\textbf{Intuition:}
The traces of a matrix tell us about the "moments" of its eigenvalues. If every single moment of the eigenvalues is zero, the only possible way for that to happen is if all the eigenvalues themselves are zero. A matrix with only zero eigenvalues is fundamentally "empty" in terms of growth, eventually shrinking to zero when multiplied repeatedly.

\par\vspace{1.5em}

\subsection*{Part (3): Same Characteristic Polynomial}

\textbf{Prerequisites:} 
Two matrices have the same characteristic polynomial if and only if they have the same set of eigenvalues (counting multiplicity).

\textbf{Steps:}
\begin{enumerate}
    \item Let $e_k(A)$ and $e_k(B)$ be the coefficients of the characteristic polynomials of $A$ and $B$.
    \item Newton's identities provide a bijective (one-to-one) mapping between the set of traces $\{p_1, \dots, p_n\}$ and the set of coefficients $\{e_1, \dots, e_n\}$ in characteristic zero.
    \item Therefore, $p_k(A) = p_k(B)$ for all $k$ implies $e_k(A) = e_k(B)$ for all $k$, and vice versa.
\end{enumerate}

\textbf{Intuition:}
This result shows that the traces of the powers of a matrix contain exactly the same amount of information as the characteristic polynomial itself. If you know the trace of $A, A^2, \dots, A^n$, you have "fingerprinted" the matrix's eigenvalues completely.
\end{solution}

\begin{exercise}
\label{algebra:2020:2}
\begin{enumerate}
\item Let $M$ be a finitely generated $R$-module and $\mathfrak{a} \subset R$ an ideal. Suppose $\phi: M \to M$ is an $R$-module map such that $\phi(M) \subseteq \mathfrak{a}M$. Prove that there is a monic polynomial $p(t) \subset R[t]$ with coefficients from $\mathfrak{a}$ such that $p(\phi) = 0$.

Hint: $p(t)$ is basically just the characteristic polynomial.

\item If $M$ is a finitely generated $R$-module such that $\mathfrak{a}M = M$ for some ideal $\mathfrak{a} \subset R$, then there exists $x \in R$ such that $1 - x \in \mathfrak{a}$ and $xM = 0$.
\end{enumerate}
\end{exercise}

\begin{solution}
\subsection*{Part (1): Generalized Cayley-Hamilton}

\textbf{Prerequisites:} 
Understanding of \textit{modules} (which are like vector spaces over a ring), \textit{generators}, and \textit{adjugate matrices}. A module is finitely generated if there's a finite set of elements $\{m_1, \dots, m_n\}$ such that every element in $M$ is a linear combination of these.

\textbf{Steps:}
\begin{enumerate}
    \item Since $M$ is generated by $m_1, \dots, m_n$ and $\phi(m_i) \in \mathfrak{a}M$, we can express each $\phi(m_i)$ as:
    \[ \phi(m_i) = \sum_{j=1}^n a_{ij} m_j, \quad \text{where each } a_{ij} \in \mathfrak{a}. \]
    \item Rearrange this into a system of equations: $\sum_{j=1}^n (\delta_{ij} \phi - a_{ij}) m_j = 0$.
    \item Think of this as a matrix equation $A \mathbf{m} = 0$, where $A$ is an $n \times n$ matrix with entries in the ring $R[\phi]$, and $\mathbf{m}$ is the vector of generators.
    \item Multiply by the adjugate matrix $\mathrm{adj}(A)$. From linear algebra, $\mathrm{adj}(A) A = \det(A) I$.
    \item This gives $\det(A) m_i = 0$ for every generator $m_i$. Since these generate $M$, we have $\det(A) M = 0$.
    \item The determinant $p(\phi) = \det(\phi I - (a_{ij}))$ is a monic polynomial in $\phi$ with coefficients in $\mathfrak{a}$.
\end{enumerate}

\textbf{Intuition:}
This is a beautiful extension of the Cayley-Hamilton theorem from linear algebra. In a vector space, a linear map satisfies its own characteristic polynomial. Here, we show that even in a more general "module" setting, if a map sends elements into a "smaller" part of the module (the ideal $\mathfrak{a}M$), the map must satisfy a polynomial whose coefficients reflect that restriction.

\par\vspace{1.5em}

\subsection*{Part (2): Nakayama's Lemma}

\textbf{Prerequisites:} 
This part is a direct application of Part (1). It is a version of the famous \textit{Nakayama's Lemma}.

\textbf{Steps:}
\begin{enumerate}
    \item Assume $\mathfrak{a}M = M$. Let $\phi$ be the identity map $\mathrm{Id}_M$.
    \item Then $\phi(M) = M = \mathfrak{a}M$. By Part (1), there is a monic polynomial $p(t) = t^n + a_{n-1} t^{n-1} + \dots + a_0$ with $a_i \in \mathfrak{a}$ such that $p(\mathrm{Id}_M) = 0$.
    \item Applying the identity map: $(\mathrm{Id}_M^n + a_{n-1} \mathrm{Id}_M^{n-1} + \dots + a_0) M = 0$.
    \item This simplifies to $(1 + a_{n-1} + \dots + a_0) M = 0$.
    \item Let $1-x = -(a_{n-1} + \dots + a_0)$. Since all $a_i \in \mathfrak{a}$, their sum is in $\mathfrak{a}$. Thus $1-x \in \mathfrak{a}$, and $x M = 0$ (or vice versa, the notation depends on how you define $x$). Specifically, if we set $1-x = \sum a_i$, then $x \cdot M = 0$.
\end{enumerate}

\textbf{Intuition:}
Nakayama's Lemma roughly says that if an ideal "consumes" a module ($\mathfrak{a}M=M$), then there is some element "close to 1" (specifically $1-x \in \mathfrak{a}$) that acts as zero on the whole module. In many cases (like in local rings), this implies that the module $M$ must actually be zero. It's a powerful tool for showing that something is zero by checking it "modulo" an ideal.
\end{solution}

\begin{exercise}
\label{algebra:2020:3}
Let $R = F[x, y]/(y^2 - x^2 - x^3)$ for some field $F$.

\begin{enumerate}
\item Prove that $R$ is an integral domain.

\item Compute the normalization of $R$ (i.e., the integral closure of $R$ in its field of fraction).
\end{enumerate}
\end{exercise}

\begin{solution}
\subsection*{Part (1): Proving $R$ is an Integral Domain}

\textbf{Prerequisites:} 
A quotient ring $F[x,y]/(f)$ is an integral domain if and only if the polynomial $f$ is \textit{irreducible}. We use the fact that in a polynomial ring, a square must have even degree.

\textbf{Steps:}
\begin{enumerate}
    \item Consider $f(x, y) = y^2 - x^2(1+x)$. This is a polynomial in $y$ of degree 2.
    \item For a quadratic $y^2 - g(x)$ to be reducible, $g(x)$ must be a perfect square in the ring $F[x]$.
    \item Here $g(x) = x^2(1+x)$. The degree of $g(x)$ is 3.
    \item In any polynomial ring, the degree of a perfect square $h(x)^2$ is $2 \cdot \deg(h)$, which is always an even number.
    \item Since 3 is odd, $x^2(1+x)$ cannot be a square. Therefore, $f(x, y)$ is irreducible, and $R$ is an integral domain.
\end{enumerate}

\textbf{Intuition:}
The equation $y^2 = x^2 + x^3$ defines a curve in the plane. Proving $R$ is an integral domain is like saying this curve doesn't "break apart" into two or more simpler curves (like two lines). Because the right-hand side isn't a square, the "square root" $y$ isn't a simple polynomial in $x$, so the relation is as simple as it can be.

\par\vspace{1.5em}

\subsection*{Part (2): Normalization}

\textbf{Prerequisites:} 
The \textit{normalization} is the process of adding all "missing" roots of monic polynomials to the ring. For curves, this often involves "smoothing out" a singularity (like the crossing point at the origin).

\textbf{Steps:}
\begin{enumerate}
    \item In the field of fractions, we can define $t = y/x$.
    \item Square this: $t^2 = y^2/x^2 = (x^2 + x^3)/x^2 = 1 + x$.
    \item From $t^2 = 1+x$, we get $x = t^2 - 1$.
    \item Substitute this into $y = xt$: $y = (t^2 - 1)t = t^3 - t$.
    \item Note that $t$ is \textit{integral} over $R$ because it satisfies the monic equation $t^2 - (1+x) = 0$, where $1+x \in R$.
    \item The ring $F[t]$ is a polynomial ring, which is always integrally closed (normal).
    \item Since $R = F[t^2-1, t^3-t] \subset F[t]$ and $t$ is in the field of fractions, the normalization of $R$ is exactly $F[t]$.
\end{enumerate}

\textbf{Intuition:}
The curve $y^2 = x^2 + x^3$ has a "node" (a self-intersection) at the origin $(0,0)$. By introducing the new variable $t = y/x$, we are essentially "unwrapping" the curve. At the origin, $y/x$ has two possible values (the slopes of the two branches), so $t$ distinguishes between them. The resulting ring $F[t]$ represents a smooth line, which is the "normalized" version of our original curvy, self-intersecting shape.
\end{solution}

\begin{exercise}
\label{algebra:2020:4}
Let $p$ and $\ell$ be two prime numbers and $[\ell_x] = 1 + x + \cdots + x^{\ell-1}$ denote the $\ell$-th cyclotomic polynomial.

\begin{enumerate}
\item Prove that $[\ell_x]$ is an irreducible element of $\mathbb{Q}[x]$.

\item Show that $[\ell_x]$ is divisible by $x - 1$ in $\mathbb{F}_p[x]$ if $p = \ell$.

\item Suppose $p \neq \ell$. Let $a$ be the order of $p$ in $\mathbb{F}_\ell$. Show that $a$ is the first value of $m$ for which the group $\mathrm{GL}_m(\mathbb{F}_p)$ of invertible $m \times m$ matrices with entries from $\mathbb{F}_p$ contains an element of order $\ell$.

Hint: Derive and use the formula for the number of elements in $\mathrm{GL}_m(\mathbb{F}_p)$.
\end{enumerate}
\end{exercise}

\begin{solution}
\subsection*{Part (1): Irreducibility of $\Phi_\ell(x)$}

\textbf{Prerequisites:} 
\textit{Eisenstein's Criterion} is a standard way to check if a polynomial is irreducible. It requires a prime $p$ such that $p$ divides all coefficients except the leading one, and $p^2$ does not divide the constant term.

\textbf{Steps:}
\begin{enumerate}
    \item Let $f(x) = \Phi_\ell(x) = \frac{x^\ell - 1}{x - 1}$. It's hard to apply Eisenstein directly here.
    \item Perform a change of variables: $g(x) = f(x+1) = \frac{(x+1)^\ell - 1}{(x+1) - 1} = \frac{(x+1)^\ell - 1}{x}$.
    \item Expand using the Binomial Theorem: $g(x) = \frac{1}{x} \left( x^\ell + \binom{\ell}{\ell-1}x^{\ell-1} + \dots + \binom{\ell}{1}x + 1 - 1 \right) = x^{\ell-1} + \binom{\ell}{\ell-1}x^{\ell-2} + \dots + \binom{\ell}{1}$.
    \item Since $\ell$ is prime, every binomial coefficient $\binom{\ell}{k}$ for $1 \leq k < \ell$ is divisible by $\ell$.
    \item The constant term is $\binom{\ell}{1} = \ell$, which is divisible by $\ell$ but \textit{not} by $\ell^2$.
    \item By Eisenstein's Criterion, $g(x)$ is irreducible, so $f(x)$ is also irreducible.
\end{enumerate}

\textbf{Intuition:}
Symmetry is key. The polynomial $\Phi_\ell(x)$ looks very uniform. By "shifting" it by 1 ($x \to x+1$), we break the symmetry in a way that reveals the prime $\ell$ lurking in the coefficients, allowing us to use a powerful divisibility test (Eisenstein).

\par\vspace{1.5em}

\subsection*{Part (2): Divisibility in Characteristic $p$}

\textbf{Prerequisites:} 
In a field of characteristic $p$, we have the "Freshman's Dream": $(a+b)^p = a^p + b^p$.

\textbf{Steps:}
\begin{enumerate}
    \item If $p = \ell$, we are working in $\mathbb{F}_\ell$.
    \item The polynomial $x^\ell - 1$ becomes $(x-1)^\ell$ because $(x-1)^\ell = x^\ell + (-1)^\ell = x^\ell - 1$ (since $\ell$ is prime, if $\ell=2$, $-1=1$; if $\ell$ is odd, $(-1)^\ell = -1$).
    \item Then $\Phi_\ell(x) = \frac{x^\ell - 1}{x - 1} = \frac{(x-1)^\ell}{x-1} = (x-1)^{\ell-1}$.
    \item This is clearly divisible by $(x-1)$, provided $\ell-1 \geq 1$ (which holds for any prime $\ell$).
\end{enumerate}

\textbf{Intuition:}
When the characteristic of the field matches the degree of the roots we are looking for, things "collapse." Instead of having $\ell$ distinct roots, all the roots of $x^\ell-1$ merge into a single root at $x=1$ with high multiplicity.

\par\vspace{1.5em}

\subsection*{Part (3): Elements of order $\ell$ in $GL_m(\mathbb{F}_p)$}

\textbf{Prerequisites:} 
The size of the general linear group is $|GL_m(\mathbb{F}_p)| = (p^m-1)(p^m-p)\dots(p^m-p^{m-1})$. \textit{Cauchy's Theorem} states that a group has an element of order $\ell$ if and only if $\ell$ divides the group's order.

\textbf{Steps:}
\begin{enumerate}
    \item Factor out powers of $p$ from the order formula: $|GL_m(\mathbb{F}_p)| = p^{m(m-1)/2} \prod_{i=1}^m (p^i - 1)$.
    \item Since $p \neq \ell$, $\ell$ can only divide $|GL_m(\mathbb{F}_p)|$ if it divides one of the factors $(p^i - 1)$.
    \item This means $p^i \equiv 1 \pmod{\ell}$ for some $1 \leq i \leq m$.
    \item The smallest such $i$ is, by definition, the \textit{order} of $p$ in the multiplicative group $(\mathbb{Z}/\ell\mathbb{Z})^\times$. Let's call this order $a$.
    \item Thus, $\ell$ divides the group order for the first time when $m = a$. By Cauchy's Theorem, this is the first time an element of order $\ell$ appears.
\end{enumerate}

\textbf{Intuition:}
In a finite field, the existence of certain "rotations" or "symmetries" (elements of order $\ell$) depends on whether the field's size allows for that specific arithmetic. The order $a$ tells us exactly how many dimensions we need to "room" an operation that cycles through $\ell$ states.
\end{solution}

\begin{exercise}
\label{algebra:2020:5}
Let $p \geq 3$ be a prime number and let $\mathbb{Z}_p$ be the ring of $p$-adic integers.

\begin{enumerate}
\item Show that an element in $1 + p\mathbb{Z}_p$ is a $p$-th power in $\mathbb{Z}_p$ if and only if it lives in $1 + p^2\mathbb{Z}_p$.

\item Let $\mathbb{Z}_p^\times$ denote the group of units in $\mathbb{Z}_p$. Show that there exist $a, b, c \in \mathbb{Z}_p^\times$ such that $a^p + b^p = c^p$ if and only if
\[\sum_{i=1}^{p-1} i^{p-2} t^i \equiv 0 \pmod{p}\]
for some integer $t \in \{2, 3, \ldots, p-1\}$. (In particular, this condition holds for $p = 7$ by taking $t = 3$. Therefore, Fermat's Last Theorem does not hold for $\mathbb{Z}_7$.)
\end{enumerate}
\end{exercise}

\begin{solution}
\subsection*{Part (1): $p$-th powers in $p$-adic integers}

\textbf{Prerequisites:} 
The \textit{$p$-adic logarithm} and \textit{exponential} functions. For $p \geq 3$, these provide an isomorphism between the multiplicative group $1 + p\mathbb{Z}_p$ and the additive group $p\mathbb{Z}_p$.

\textbf{Steps:}
\begin{enumerate}
    \item Let $x \in 1 + p\mathbb{Z}_p$. Then $x$ is a $p$-th power if $x = y^p$ for some $y$.
    \item Applying the $p$-adic logarithm: $\log(x) = \log(y^p) = p \log(y)$.
    \item This means $\log(x)$ must be divisible by $p$ in the additive space $p\mathbb{Z}_p$.
    \item Thus $\log(x) \in p(p\mathbb{Z}_p) = p^2 \mathbb{Z}_p$.
    \item Since $x \approx 1 + \log(x)$, the condition $\log(x) \in p^2 \mathbb{Z}_p$ is equivalent to $x \in 1 + p^2 \mathbb{Z}_p$.
\end{enumerate}

\textbf{Intuition:}
The $p$-adic world behaves like a "stretched" version of regular arithmetic. Taking a $p$-th power "pushes" an element further away from 1 by exactly one power of $p$. If you start at $1+p$, you end up at $1+p^2$. This makes $p$-th powers easy to spot just by looking at their "depth" or precision.

\par\vspace{1.5em}

\subsection*{Part (2): Fermat's "Not" Last Theorem in $\mathbb{Z}_p$}

\textbf{Prerequisites:} 
\textit{Hensel's Lemma} (which lets us lift solutions from $\pmod p$ to $\mathbb{Z}_p$) and the \textit{Teichmüller lift} (a way to represent elements of $\mathbb{F}_p$ as actual $p$-adic numbers).

\textbf{Steps:}
\begin{enumerate}
    \item If $a^p + b^p = c^p$ in $\mathbb{Z}_p$, then modulo $p^2$, we have $a^p + b^p \equiv c^p \pmod{p^2}$.
    \item By Fermat's Little Theorem, $x^p \equiv x \pmod p$. More precisely, $x^p \equiv \omega(x)^p \pmod{p^2}$ where $\omega(x)$ is the Teichmüller lift.
    \item Setting $t = b/a \pmod p$, the condition becomes $1 + t^p \equiv (1+t)^p \pmod{p^2}$.
    \item Expanding the right side: $(1+t)^p = 1 + t^p + \sum_{i=1}^{p-1} \binom{p}{i} t^i$.
    \item We need $\sum_{i=1}^{p-1} \binom{p}{i} t^i \equiv 0 \pmod{p^2}$, which means $\sum_{i=1}^{p-1} \frac{1}{p}\binom{p}{i} t^i \equiv 0 \pmod p$.
    \item Using the identity $\frac{1}{p}\binom{p}{i} \equiv \frac{(-1)^{i-1}}{i} \pmod p$ and $i^{p-2} \equiv 1/i \pmod p$, we get the polynomial condition in the problem.
\end{enumerate}

\textbf{Intuition:}
Fermat's Last Theorem is famous because it's impossible for integers. But $p$-adic numbers are much more "flexible"—they allow for a lot of room to "wiggle" a solution into existence. The condition given is the "gatekeeper" that determines if that wiggle room is enough to satisfy the equation at the lowest level of precision.
\end{solution}

\begin{exercise}
\label{algebra:2020:6}
Recall that a metric space is called spherically complete if any decreasing sequence of closed balls has nonempty intersection.

Let $p$ be a prime number and let $\mathbb{Q}_p$ be the field of $p$-adic numbers. For every integer $n \geq 1$, consider the finite extension $\mathbb{Q}_p(\mu_{p^n})$ of $\mathbb{Q}_p$ generated by all $p^n$-th roots of unity. Let $\mathbb{Q}_p(\mu_{p^\infty}) = \cup_{n \geq 1} \mathbb{Q}_p(\mu_{p^n})$.

Question: Which of the following are spherically complete? Explain why.

\begin{enumerate}
\item $\mathbb{Q}_p$.

\item $\mathbb{Q}_p(\mu_{p^n})$.

\item $\mathbb{Q}_p(\mu_{p^\infty})$.

\item $\overline{\mathbb{Q}_p(\mu_{p^\infty})}$, the completion of $\mathbb{Q}_p(\mu_{p^\infty})$.
\end{enumerate}
\end{exercise}

\begin{solution}
\subsection*{General Rule for Spherical Completeness}

\textbf{Prerequisites:} 
In a standard metric space like $\mathbb{R}$, "complete" means every Cauchy sequence converges. In the "weird" world of $p$-adic numbers (non-archimedean), we have a stronger property: \textit{spherical completeness}. This means that even if the radii of a sequence of nested balls don't go to zero, they must still have a common point.

\textbf{Steps \& Analysis:}
\begin{enumerate}
    \item \textbf{$\mathbb{Q}_p$ is spherically complete.} It is a locally compact field with a discrete valuation. For discrete fields, completeness and spherical completeness are the same thing. Since $\mathbb{Q}_p$ is the completion of $\mathbb{Q}$, it fits the bill.
    
    \item \textbf{$\mathbb{Q}_p(\mu_{p^n})$ is spherically complete.} This is a finite extension of $\mathbb{Q}_p$. Finite extensions preserve the "discrete" nature of the valuation and the property of being complete.
    
    \item \textbf{$\mathbb{Q}_p(\mu_{p^\infty})$ is NOT spherically complete.} This is an infinite union of fields. It's like taking all fractions with power-of-2 denominators—it's full of "holes" that haven't been filled yet. It isn't even complete in the basic sense.
    
    \item \textbf{$\overline{\mathbb{Q}_p(\mu_{p^\infty})}$ is NOT spherically complete.} This is the trickiest one. Although it is "complete" (no holes for sequences), its "value group" (the possible sizes of elements) is not discrete—it's "dense." In such fields, you can have a sequence of balls that shrink but never quite reach a point.
\end{enumerate}

\textbf{Intuition:}
Imagine a set of nested Russian dolls. If the dolls get infinitely thin (radii go to zero), you'd expect to find a "tiny center" (a point). But if the dolls stay "fat" (radii go to, say, 0.5), spherical completeness guarantees there's still something inside. $\mathbb{Q}_p$ is "solid" enough to have this, but infinite extensions like $\mathbb{C}_p$ are too "porous" at the microscopic level.
\end{solution}
